\section{Фабрика: семь шин, не кроссбар}
\label{sec:fabric}

\subsection{Интерконнект --- это семь прямых}

Двадцать одна парная связь среди семи секторов потребовала бы на обычной
фабрике сети из $\binom{7}{2}=21$ ребра или кроссбара $7\times7$. Коллапс
(Теорема~\ref{thm:collapse}, L1) заменяет это семью прямыми $\mathsf{N}$ как
\textbf{семью тройственными шинами}. Всякая пара секторов делит ровно одну шину
(Штейнер $S(2,3,7)$); всякий сектор сидит ровно на трёх шинах; любые две шины
пересекаются ровно в одном секторе. Это не проектные цели --- это теоремы
проективной плоскости, и они дают архитектуре три свойства бесплатно:

\begin{itemize}[leftmargin=1.5em,itemsep=1pt]
  \item \textbf{Равномерная нагрузка, по симметрии.} Каждый сектор на ровно трёх
        шинах, потому что $\mathrm{PSL}(2,7)$ действует 2-транзитивно на семи
        точках; балансировка нагрузки --- следствие группы автоморфизмов, не
        планировщика.
  \item \textbf{Гарантированное рандеву, по геометрии.} Любые две шины
        встречаются в секторе, поэтому любые два сектора, которым надо обменяться
        состоянием, имеют детерминированную точку встречи --- нет недостижимой
        пары и нет поиска маршрута.
  \item \textbf{Детерминированная починка, по закону замыкания.} Если прямая
        шина для $(i,j)$ деградировала, состояние регенерируется через
        \emph{медиатора} $k^\ast(i,j)$ --- третий сектор их прямой --- и только
        через него. Резервная маршрутизация не нуждается в таблицах: следующий
        хоп вычисляется из пары.
\end{itemize}

\subsection{Маршрутизация третьего порядка и преимущество интеграции $9\times$}

Обычный маршрутизатор двигает данные парно и, в силу слепого пятна первого
порядка (\S\ref{sec:collapse}), неспособен видеть структуру, которую несёт:
всякая равновзвешенная парная фабрика --- функция одного лишь $J-I$ и не может
отличить проводку Фано от бесструктурной all-to-all. \TALOS{} маршрутизирует на
\emph{тройках}: передача --- взаимодействие $(i,j\mid k^\ast)$, опосредованное
третьим сектором прямой. Это аппаратная форма принципа третьего порядка, и она
несёт измеримую цену за ошибку.

\begin{claimbox}[Преимущество линие-осознанной маршрутизации \status{Т} арифметика]
Когерентность, толкаемая через секторы, не делящие прямую, ослабляется точной
константой контракции канала $1/9$ за грубый хоп ($\Phi\to\Phi/9$, счёт
инцидентностей Штейнера: населённость выживает во всех трёх своих прямых,
$3\times\tfrac13=1$, но когерентность выживает лишь в одной из трёх, $\times\tfrac13$,
в квадрате $=\tfrac19$). Маршрутизатор, слепой к линейной структуре, поэтому
теряет фактор \textbf{девять} в интеграции на хоп; линие-осознанный
маршрутизатор --- удерживая всякое взаимодействие на прямой его медиатора ---
сохраняет её. Интеграция на операцию в $9\times$ выше, чем у линие-слепой фабрики.
\end{claimbox}

Это резчайшая причина, почему топология не может быть запоздалой мыслью: общий,
линие-агностичный mixture-of-experts маршрутизатор (умолчание в больших нейронных
ускорителях) платит $\times\tfrac19$ на каждом межэкспертном хопе, и никакая
полоса это не возвращает, потому что потеря --- в \emph{интеграции}, а не в
пропускной способности.

\begin{trybox}[Попробуй сам --- проложи маршрут одним пальцем]
Поставь палец на фигуру Фано (Рис.~\ref{fig:collapse}, справа). Игра:
донеси сообщение от $S$ к $E$. Правило: двигаться можно только вдоль линии,
а пара встречается \emph{только через третью точку собственной линии}.
Найди линию, содержащую и $S$, и $E$, --- это $\{S,E,U\}$, --- значит,
сообщение идёт $S\to U\to E$: медиатор --- $U$, и искать было нечего,
потому что пара сама \emph{называет} своё реле. Теперь порви прямой путь от
$A$ к $O$ и переложи маршрут: их линия --- $\{A,O,U\}$, снова через $U$ ---
починка без таблиц маршрутизации (\S\ref{sec:ecc}). Все $21$ пара играют в
ту же одноходовую игру; полная таблица медиаторов --- в
Приложении~\ref{app:fano}.
\end{trybox}

\begin{table}[H]
\centering\small
\caption{Интерконнект: кроссбар $7\times7$ против семишинной фабрики Фано.
Шинная фабрика есть также код и инструкция (Теорема~\ref{thm:collapse}); кроссбар
--- ни то, ни другое.}
\label{tab:crossbar-vs-fano}
\begin{tabular}{@{}p{3.6cm}p{3.6cm}p{4.5cm}@{}}
\toprule
\textbf{Свойство} & \textbf{Кроссбар $7\times7$} & \textbf{Семишинная Фано} \\
\midrule
Связи & 21 парная / полный кроссбар & 7 тройственных шин \\
Нагрузка на узел & зависит от планировщика & ровно 3 (по $\mathrm{PSL}(2,7)$) \\
Достижимость & all-to-all, поиском & любая пара делит одну шину \\
Reroute при отказе & таблицы маршрутизации & медиатор $k^\ast(i,j)$, без таблиц \\
Видит свою структуру? & нет (слепа к 1-му порядку) & да (третий порядок) \\
Также ECC? & нет & да (Хэмминг $[7,4,3]$, \S\ref{sec:ecc}) \\
Также алгебра ISA? & нет & да (октонионные $f_{ijk}$) \\
\bottomrule
\end{tabular}
\end{table}
